The following is a significant improvement due to a discussion with Ingo Blechschmidt and Marc Nieper-Wißkirchen, on the original observation that formally étale surjective maps between smooth irreducible separated schemes always have a section.

\begin{lemma}
  Let $X$ be an étale-overt irreducible scheme whose identity types are étale-sheaves, $Y$ a connected scheme such that $\|Y\|_\et$ and $f:X\to Y$ unramified separated
  \footnote{i.e. all fibers are unramified separated schemes -- we could just say map with decidable fibers here as well and that $Y$ just needs to be a set which is connected in the sense that maps to Bool are constant.}.
  Then any two sections of $f$ are equal.
\end{lemma}

\begin{proof}
  Let $s,t:Y\to X$ be sections of $f$.
  The propositions $S(x)\colonequiv s(f(x))=x$ and $T(x)\colonequiv t(f(x))=x$
  are equivalent to $x\in\im(s)$ and $x\in\im(t)$ for all $x:X$.
  As $s(f(x))$ and $x$ are both in $\fib_f(f(x))$,
  $S(x)$ is equivalent to an identity in this fiber and therefore decidable,
  since the fibers of $f$ are unramified separated schemes.
  The same holds for $T(x)$ which means that $\im(s)$ and $\im(t)$ are open in $X$ (and closed).
  By irreducibility and étale-overtness, we have $\|\im(s)\cap\im(t)\|_\et$.

  To show $s=t$ which is an étale-sheaf and a proposition,
  we can assume $u:Y$ with $s(u)=t(u)$.
  For any $y:Y$, the proposition $s(y)=t(y)$ is decidable.
  Since $Y$ is connected, $s(y)=t(y)$ has to be constant on all of $Y$ and we know $s(u)=t(u)$.
\end{proof}

\begin{proposition}
  Let $X$ be irreducible étale-overt scheme and $Y$ a scheme.
  Any unramified separated surjective map $f:X\to Y$ is an equivalence.
\end{proposition}

\begin{proof}
  By Zariski-local choice, there are local sections which glue to a unique global section $s:Y\to X$ using the lemma.
  As in the lemma, the image of $s$ is an open $U\subseteq X$, which is also closed.
  Assume $x\notin U$, then we know it is different from $x^\prime\colonequiv s(f(x))$ and lies in the open complement of $U$, but by irreducibility, this complement has to be empty.
  By double-negation stability of opens, $s$ is therefore surjective.
  This means $s$ is an equivalence and therefore its one-sided inverse $f$.
\end{proof}

A possibly related statement:

\begin{proposition}
  \label{cohom-discrete-coeff}
  Let $X$ be an irreducible, étale-overt, pointed scheme.
  If $G$ is a group which is an étale sheaf with decidable equality, then:
  \[
    H^1(X,G)=0.
  \]
\end{proposition}

\begin{proof}
  Let $\chi:X\to K(G,1)$.
  By Zariski-local choice, we have an open cover $U_1,\dots,U_n$ and
  trivializations $s_i:(x:U_i)\to \chi_x=\ast$.
  Let $k$ such that $\ast:U_k$.
  Then for any $x:X$ and $i$ such that $x:U_i$,
  we have \(U_k\cap U_i\neq \emptyset\) and $\|U_k\cap U_i\|_\et$.
  
  Let \(D_i\colonequiv (g:G)\times ((u:U_i\cap U_k)\to gs_i(u)=s_k(u))\), still assuming \(x:U_i\).
  \(D_i\) is an étale sheaf and using that equality in \(G\) is double negation stable and
  \(U_k\cap U_i\neq \emptyset\), we have that \(D_i\) is a proposition.
  We can therefore use \(y:U_k\cap U_i\) to show that \(D_i\) is inhabited,
  but then \(g_i\) is just given as the difference of \(s_k(y)\) and \(s_i(y)\),
  which are both elements of the same \(G\)-torsor.

  Let \(\widetilde{s}_i\colonequiv g_is_i\).
  We conclude by showing that \(\widetilde{s}_i(x)=\widetilde{s}_j(x)\) holds for \(x:U_i\cap U_j\).
  The goal is a propositional étale sheaf and \(U_k\cap U_i\cap U_j\) is non-empty,
  so we can assume \(y:U_k\cap U_i\cap U_j\).
  But then we have \(s_k(y)=\widetilde{s}_i(y)=\widetilde{s}_j(y)\).
  Now \(U_i\cap U_j\) is connected by being irreducible inhabited, so \(\widetilde{s}_i(u)=\widetilde{s}_j(u)\) holds for all \(u:U_i\cap U_j\).
\end{proof}

\begin{example}
  \(R^n,\bP^n\) and any irreducible pointed smooth scheme can be used as ``\(X\)'' in \Cref{cohom-discrete-coeff}. Examples for the group ``\(G\)'' include \(\Z\), any finite group and if we have \(\ell\neq 0\) in \(R\), then \(\mu_\ell\colonequiv \{x:R\mid x^\ell=1\}\) is an étale group scheme and therefore an example as well.
\end{example}
